\section{Introduzione}
\subsection{Scopo del documento}

Il presente documento ha l'obiettivo di definire e descrivere l'architettura adottata per lo sviluppo$^G$ del sistema SmartOrder.\\ 
Il documento delinea le scelte architetturali, tecnologiche e di design adottate dal gruppo per la realizzazione del prodotto.\\
In particolare il documento si propone di definire: 
\begin{itemize}[itemsep=0pt, parsep=1pt, label=$\scriptstyle\bullet$]
    \item L'architettura del sistema, comprensiva delle componenti principali e delle loro interazioni;
    \item Le tecnologie e i framework$^G$ selezionati per lo sviluppo$^G$, con particolare attenzione alle motivazioni tecniche che ne hanno guidato l’adozione;
    \item I design pattern$^G$ applicati e le ragioni tecnologiche che ne hanno guidato l'adozione;
    \item Il tracciamento dei requisiti, che attesta la copertura di quanto stabilito nel documento di Analisi dei Requisiti.
\end{itemize}
Il documento è soggetto ad aggiornamenti durante tutto il ciclo di vita del progetto$^G$, al fine di recepire nuove decisioni architetturali o eventuali modifiche legate all’evoluzione del progetto.

\subsection{Scopo del progetto}
SmartOrder è un prodotto ideato dall'azienda Ergon Informatica per automatizzare il processo$^G$ di gestione degli ordini di acquisto aziendali tramite l'analisi di input multimodali quali testo, immagini o audio. 
Attraverso l'utilizzo del Natural Language Processing$^G$ (NLP$^G$), SmartOrder interpreta queste richieste spesso non strutturate, estraendo le informazioni rilevanti quali prodotti e quantità, trasformandole in ordini strutturati pronti per l'integrazione nei sistemi ERP$^G$ aziendali.

\subsection{Glossario}
Per tutti i termini tecnici, gli acronimi e le definizioni utilizzate nel documento si rimanda al \underline{\href{https://nullpointersgroup.github.io/Documentazione/output/PB/Documenti\%20Interni/Glossario.pdf}{Glossario$^G$}}, disponibile come documento separato.\\
Ogni parola presente nel Glossario$^G$ viene segnata come segue:
\begin{center}
	termine$^{G}$
\end{center}


\subsection{Riferimenti}
\subsubsection{Riferimenti normativi}

\begin{itemize}[itemsep=5pt, parsep=5pt, label=$\scriptstyle\bullet$]

	\item \textbf{Norme di Progetto$^G$, versione 2.0.0}\\
	\url{https://nullpointersgroup.github.io/Documentazione/output/PB/Documenti%20Interni/Norme_di_Progetto.pdf}\\[3pt]
	\textbf{Ultima consultazione: 25 marzo 2026}
	
	\item \textbf{Capitolato$^G$ C8 - Ergon Informatica Srl - SmartOrder}\\
	\url{https://www.math.unipd.it/~tullio/IS-1/2025/Progetto/C8.pdf}\\[3pt]
	\textbf{Ultima consultazione: 10 aprile 2026}

\end{itemize}

\subsubsection{Riferimenti informativi}
\begin{itemize}[itemsep=5pt, parsep=5pt, label=$\scriptstyle\bullet$]

	\item \textbf{Glossario$^G$, versione 2.0.0}\\
	 \url{https://nullpointersgroup.github.io/Documentazione/output/PB/Documenti%20Interni/Glossario.pdf}\\[3pt]
	\textbf{Ultima consultazione: 11 marzo 2026}
	
	\item \textbf{Analisi dei requisiti$^G$, versione 2.0.0}\\
	 \url{https://nullpointersgroup.github.io/Documentazione/output/PB/Documenti%20Interni/Analisi_dei_Requisiti.pdf}\\[3pt]
	\textbf{Ultima consultazione: 07 aprile 2026}
	
	\item \textbf{Software Architecture Patterns, SWE 2025-2026}\\
	 \url{https://www.math.unipd.it/~rcardin/swea/2022/Software%20Architecture%20Patterns.pdf}\\[3pt]
	\textbf{Ultima consultazione: 11 marzo 2026}
	
	\item \textbf{Dependency Injection, SWE 2025-2026}\\
	 \url{https://www.math.unipd.it/~rcardin/swea/2022/Design%20Pattern%20Architetturali%20-%20Dependency%20Injection.pdf}\\[3pt]
	\textbf{Ultima consultazione: 11 marzo 2026}
	
	\item \textbf{Design Pattern$^G$ Creazionali, SWE 2025-2026}\\
	 \url{https://www.math.unipd.it/~rcardin/swea/2022/Design%20Pattern%20Creazionali.pdf}\\[3pt]
	\textbf{Ultima consultazione: 11 marzo 2026}
	
	\item \textbf{Design Pattern$^G$ Strutturali, SWE 2025-2026}\\
	 \url{https://www.math.unipd.it/~rcardin/swea/2022/Design%20Pattern%20Strutturali.pdf}\\[3pt]
	\textbf{Ultima consultazione: 11 marzo 2026}

\end{itemize}